\documentclass{article}
\usepackage[margin=28mm]{geometry}
\usepackage{amsmath,amssymb}
\usepackage{xurl}
\usepackage[hidelinks]{hyperref}
% Shared commands for notes in docs/paper-gaps/.

\DeclareUnicodeCharacter{2080}{\ifmmode{}_{0}\else\textsubscript{0}\fi}
\DeclareUnicodeCharacter{2081}{\ifmmode{}_{1}\else\textsubscript{1}\fi}
\DeclareUnicodeCharacter{2082}{\ifmmode{}_{2}\else\textsubscript{2}\fi}
\DeclareUnicodeCharacter{2099}{\ifmmode{}_{n}\else\textsubscript{n}\fi}

\newcommand{\Lean}{\textsc{Lean}}
\ExplSyntaxOn
\NewDocumentCommand{\leanid}{m}{
  \ifmmode
    \textsf{\detokenize{#1}}
  \else
    \group_begin:
    \urlstyle{sf}
    \tl_set:Nn \l_tmpa_tl {#1}
    \tl_replace_all:Nnn \l_tmpa_tl {\_} {_}
    \exp_args:NV \path \l_tmpa_tl
    \group_end:
  \fi
}
\ExplSyntaxOff
\providecommand{\tr}{\operatorname{tr}}
\newcommand{\tnleanIssueBase}{https://github.com/LionSR/TNLean/issues/}
\newcommand{\tnleanPrBase}{https://github.com/LionSR/TNLean/pull/}
\newcommand{\tnleanDocBase}{https://sirui-lu.com/TNLean/docs/}
\newcommand{\ghissue}[1]{\href{\tnleanIssueBase#1}{GitHub issue \##1}}
\newcommand{\ghpr}[1]{\href{\tnleanPrBase#1}{PR \##1}}
\newcommand{\doclink}[1]{\href{\tnleanDocBase#1.html}{\path{#1}}}

% Dirac notation and the identity operator, as in blueprint/src/macros/common.tex.
% \providecommand keeps note-local definitions authoritative.
\usepackage{dsfont}
\providecommand{\ket}[1]{|#1\rangle}
\providecommand{\bra}[1]{\langle #1|}
\providecommand{\braket}[2]{\langle #1 | #2 \rangle}
\providecommand{\ketbra}[2]{|#1\rangle\!\langle #2|}
\providecommand{\Id}{\mathds{1}}
\providecommand{\one}{\mathds{1}}
\providecommand{\C}{\mathbb{C}}
\providecommand{\MN}[1]{M_{#1}(\C)}
\providecommand{\MD}{M_D(\C)}

% Verdict marker: \gapnote{<kind>}{<status>}. Kind is the policy
% classification (clarification, local-correction, scope-restriction,
% unfaithful, false-source, open-gap); status is open, wip (elimination
% underway), resolved, or historical. Machine-read by the site index;
% prints a verdict line.
\newcommand{\gapnote}[2]{\par\noindent\textbf{Verdict:} #1 (#2).\par}

\providecommand{\gapnote}[2]{\par\noindent\textbf{Verdict:} #1 (#2).\par}
\providecommand{\leanid}[1]{\textsf{\detokenize{#1}}}
\begin{document}
\begin{center}
    {\Large Normalization in the LVV error-reduction step}
\end{center}
\gapnote{local-correction}{resolved}

\section{The source statement}
Landau, Vazirani and Vidick~\cite{Landau2015Polynomial},
\href{https://arxiv.org/abs/1307.5143v1}{arXiv:1307.5143v1}, source
\path{algo53.tex}, line 589, give the error-reduction argument as follows.
For a witness $v$ of a $(i,\tfrac12)$-viable set they write
$\bigl|\langle\Gamma, Av/\lVert Av\rVert\rangle\bigr| \ge 1-\delta$ and
$\lVert K-A\rVert\le\delta$, with $\delta = c_\varepsilon/(2q(n))$, and conclude
$\bigl|\langle\Gamma, Kv/\lVert Kv\rVert\rangle\bigr| \ge 1-2\delta$.

\section{The point at issue}
The two displayed bounds do not imply the conclusion without a lower bound on
$\lVert Av\rVert$. In $\mathbb{C}^2$ take $\Gamma=v=e_1$, $A=\delta\,\mathbf{1}$
and $K=A+\delta\,\lvert e_2\rangle\langle e_1\rvert$. Then
$\langle\Gamma, Av\rangle/\lVert Av\rVert = 1$ and $\lVert K-A\rVert=\delta$, while
$Kv=\delta(e_1+e_2)$ has overlap $1/\sqrt2$ with $\Gamma$, which is smaller than
$1-2\delta$ for $0<\delta<1/8$.

\section{The correction}
The operator $A$ of the source is Hermitian and fixes $\Gamma$ (line 602), so
$\lVert Av\rVert \ge \lvert\langle\Gamma, Av\rangle\rvert = \lvert\langle\Gamma,
    v\rangle\rvert \ge \tfrac12$ for every witness $v$. For a unit vector $v$ with
$\lVert Av\rVert\ge c>0$, $\lvert\langle\Gamma,Av\rangle\rvert\ge(1-\delta)\lVert
    Av\rVert$ and $\lVert K-A\rVert\le\eta$, the triangle inequality gives
\begin{align}
    \lvert\langle\Gamma, Kv\rangle\rvert \ge \Bigl(1-\delta-\frac{2\eta}{c}\Bigr)
    \lVert Kv\rVert .
    \label{eq:lvv_error_reduction_overlap}
\end{align}
With $c=\tfrac12$ the error of the new viable set is $\delta+4\eta$ in place of
$2\delta$. Since the claim is stated for an arbitrary polynomial $q(n)$, the
change of constant is absorbed by applying Corollary \texttt{c:agsp} with a
larger polynomial denominator; the claim holds with the original target error.
For example, a target error $\tau>0$ is obtained by choosing the normalized
filter error at most $\tau/2$ and the sampling error at most $\tau/8$.
The vector $Kv$ is nonzero when $\eta<c$; dividing
(\ref{eq:lvv_error_reduction_overlap}) by its norm then gives a unit viable witness.

There is also a direct form that needs no separate normalized-filter estimate.
Let $P=\lvert\Gamma\rangle\langle\Gamma\rvert$ and suppose
$\lVert K-P\rVert\le\eta<1/2$. Every unit witness $v$ of a half-viable set
satisfies $\lVert Pv\rVert=|\langle\Gamma,v\rangle|\ge1/2$, and
$|\langle\Gamma,Pv\rangle|=\lVert Pv\rVert$. Applying
(\ref{eq:lvv_error_reduction_overlap}) with $A=P$ and $\delta=0$ gives
overlap at least $1-4\eta$. A power-filter estimate
$\lVert A-P\rVert\le\rho$ and a sampling estimate
$\lVert K-A\rVert\le\eta$ supply the direct hypothesis with error $\rho+\eta$.

\section{Formalization}
The perturbation estimate and its direct viable-set consequence are proved by
the following declarations.
\begin{itemize}
    \footnotesize
    \item \leanid{TNLean.MPS.AreaLaw.mul_norm_le_norm_inner_of_norm_sub_le}
    \item \leanid{TNLean.MPS.AreaLaw.IsViable.errorReduction_of_norm_sub_groundProjection_le}
\end{itemize}
\bibliographystyle{unsrt}
\bibliography{references}
\end{document}
